\textbf{Phần 1. Trắc nghiệm lựa chọn}
\setcounterpageref{ex}{0}

\begin{ex} %[2D8N37]
Tổng các số hạng bằng nhau $2 + 2 + 2 + 2$ được viết thành phép nhân là
\choice
{$2 \times 3$}
{$4 \times 4$}
{$2 \times 5$}
{\True $2 \times 4$}
\loigiai{
Số 2 được lấy 4 lần nên viết là $2 \times 4$.\\
\textbf{Đáp án của câu hỏi là D}
}
\end{ex}

\begin{ex} %[2D8N37]
Tổng $5 + 5 + 5$ viết thành phép nhân là
\choice
{\True $5 \times 3$}
{$5 \times 2$}
{$3 \times 5$}
{$5 \times 4$}
\loigiai{
$5$ lấy 3 lần là $5 \times 3$.\\
\textbf{Đáp án của câu hỏi là A}
}
\end{ex}

\begin{ex} %[2D8H37]
Phép nhân $3 \times 4$ được chuyển thành tổng các số hạng bằng nhau là
\choice
{$4 + 4 + 4$}
{\True $3 + 3 + 3 + 3$}
{$3 + 4$}
{$3 \times 3 \times 3 \times 3$}
\loigiai{
$3 \times 4 = 3 + 3 + 3 + 3$.\\
\textbf{Đáp án của câu hỏi là B}
}
\end{ex}

\begin{ex} %[2D8H37]
Tính kết quả phép nhân: $2 \times 5 = \dots$
\choice
{\True $10$}
{$8$}
{$12$}
{$15$}
\loigiai{
$2 + 2 + 2 + 2 + 2 = 10$.\\
\textbf{Đáp án của câu hỏi là A}
}
\end{ex}

\begin{ex} %[2D8H37]
Tính kết quả phép nhân: $5 \times 4 = \dots$
\choice
{$15$}
{$25$}
{\True $20$}
{$30$}
\loigiai{
$5 + 5 + 5 + 5 = 20$.\\
\textbf{Đáp án của câu hỏi là C}
}
\end{ex}

\begin{ex} %[2D8H37]
Mỗi bạn có 2 cái bút chì. Hỏi 6 bạn có tất cả bao nhiêu cái bút chì?
\choice
{$10$ cái}
{\True $12$ cái}
{$14$ cái}
{$8$ cái}
\loigiai{
$2 \times 6 = 12$ cái bút chì.\\
\textbf{Đáp án của câu hỏi là B}
}
\end{ex}

\begin{ex} %[2D8V37]
Mỗi đĩa có 5 quả cam. Hỏi 4 đĩa như thế có tất cả bao nhiêu quả cam?
\choice
{$15$ quả}
{$25$ quả}
{$10$ quả}
{\True $20$ quả}
\loigiai{
$5 \times 4 = 20$ quả cam.\\
\textbf{Đáp án của câu hỏi là D}
}
\end{ex}

\begin{ex} %[2D8V37]
Mỗi con gà có 2 cái chân. Hỏi 8 con gà có tất cả bao nhiêu cái chân?
\choice
{\True $16$ cái}
{$14$ cái}
{$18$ cái}
{$10$ cái}
\loigiai{
$2 \times 8 = 16$ cái chân.\\
\textbf{Đáp án của câu hỏi là A}
}
\end{ex}

\begin{ex} %[2D8V37]
Mỗi chuồng có 5 con thỏ. Hỏi 6 chuồng như thế có tất cả bao nhiêu con thỏ?
\choice
{$25$ con}
{$35$ con}
{\True $30$ con}
{$40$ con}
\loigiai{
$5 \times 6 = 30$ con thỏ.\\
\textbf{Đáp án của câu hỏi là C}
}
\end{ex}

\begin{ex} %[2D8V37]
Điền dấu thích hợp: $2 \times 7 \dots 5 \times 3$.
\choice
{$>$}
{$=$}
{Không so sánh được}
{\True $<$}
\loigiai{
$2 \times 7 = 14$ và $5 \times 3 = 15$. Vì $14 < 15$ nên điền dấu $<$.\\
\textbf{Đáp án của câu hỏi là D}
}
\end{ex}

\begin{ex} %[2D8V37]
Tính giá trị biểu thức: $5 \times 5 + 15$.
\choice
{$35$}
{$45$}
{\True $40$}
{$30$}
\loigiai{
$25 + 15 = 40$.\\
\textbf{Đáp án của câu hỏi là C}
}
\end{ex}

\begin{ex} %[2D8C37]
Có 4 gói kẹo, mỗi gói có 5 cái kẹo. Bé An ăn hết 3 cái kẹo. Hỏi còn lại bao nhiêu cái kẹo?
\choice
{$16$ cái}
{\True $17$ cái}
{$18$ cái}
{$15$ cái}
\loigiai{
Tổng số kẹo: $5 \times 4 = 20$ cái.\\
Còn lại: $20 - 3 = 17$ cái.\\
\textbf{Đáp án của câu hỏi là B}
}
\end{ex}

\textbf{Phần 2. Trắc nghiệm đúng sai}
\setcounterpageref{ex}{0}

\begin{ex} %[2D8H37]
Xét tính đúng sai về phép nhân:
\choiceTF
{\True Phép nhân là cách viết gọn của tổng các số hạng bằng nhau.}
{\True $2 + 2 + 2 = 2 \times 3 = 6$.}
{$5 + 5 + 5 + 5 = 5 \times 5$.}
{\True $5 + 5 + 5 + 5 = 5 \times 4 = 20$.}
\loigiai{
\begin{itemize}
    \item Phép nhân là tổng các số hạng bằng nhau: Đúng.
    \item $2 \times 3 = 6$: Đúng.
    \item $5$ lấy 4 lần là $5 \times 4$ (không phải $5 \times 5$). Ý 3 Sai, Ý 4 Đúng.
\end{itemize}
\begin{itemize}
    \item a) \textbf{Đúng}.
    \item b) \textbf{Đúng}.
    \item c) \textbf{Sai}.
    \item d) \textbf{Đúng}.
\end{itemize}
}
\end{ex}

\begin{ex} %[2D8V37]
Một bàn học có 2 bạn ngồi. Xét tính đúng sai:
\choiceTF
{\True 5 bàn học như thế có tất cả 10 bạn ($2 \times 5 = 10$).}
{\True 8 bàn học như thế có tất cả 16 bạn ($2 \times 8 = 16$).}
{10 bàn học có tất cả 18 bạn.}
{\True 10 bàn học có tất cả $2 \times 10 = 20$ bạn.}
\loigiai{
\begin{itemize}
    \item 5 bàn: $2 \times 5 = 10$: Đúng.
    \item 8 bàn: $2 \times 8 = 16$: Đúng.
    \item 10 bàn: $2 \times 10 = 20 \neq 18$. Ý 3 Sai, Ý 4 Đúng.
\end{itemize}
\begin{itemize}
    \item a) \textbf{Đúng}.
    \item b) \textbf{Đúng}.
    \item c) \textbf{Sai}.
    \item d) \textbf{Đúng}.
\end{itemize}
}
\end{ex}

\textbf{Phần 3. Trả lời ngắn}
\setcounterpageref{ex}{0}

\begin{ex} %[2D8H37]
Tính kết quả phép nhân: $2 \times 7$.
\shortans{14}
\loigiai{
\textbf{Đáp số:} $14$\\
$2 \times 7 = 14$.
}
\end{ex}

\begin{ex} %[2D8H37]
Tính kết quả phép nhân: $5 \times 6$.
\shortans{30}
\loigiai{
\textbf{Đáp số:} $30$\\
$5 \times 6 = 30$.
}
\end{ex}

\begin{ex} %[2D8V37]
Mỗi ô tô có 4 bánh xe. Hỏi 5 ô tô như thế có tất cả bao nhiêu bánh xe?
\shortans{20}
\loigiai{
\textbf{Đáp số:} $20$\\
$4 \times 5 = 20$ bánh xe.
}
\end{ex}

\begin{ex} %[2D8V37]
Mỗi hàng có 5 cây. Hỏi 7 hàng có tất cả bao nhiêu cây?
\shortans{35}
\loigiai{
\textbf{Đáp số:} $35$\\
$5 \times 7 = 35$ cây.
}
\end{ex}

\begin{ex} %[2D8V37]
Tính giá trị biểu thức: $2 \times 9 + 12$.
\shortans{30}
\loigiai{
\textbf{Đáp số:} $30$\\
$18 + 12 = 30$.
}
\end{ex}

\begin{ex} %[2D8C37]
Mỗi gói bánh có 5 cái bánh. Mẹ mua 8 gói bánh và cho các bạn 15 cái bánh. Hỏi mẹ còn lại bao nhiêu cái bánh?
\shortans{25}
\loigiai{
\textbf{Đáp số:} $25$\\
Tổng: $5 \times 8 = 40$ cái. Còn lại: $40 - 15 = 25$ cái bánh.
}
\end{ex}
